\documentclass[11pt]{article}

% Shared notes settings for Elements of AIML lectures (UPES).
% Keep this file free of \documentclass to allow reuse via \input.

\usepackage[utf8]{inputenc}
\usepackage[T1]{fontenc}
\usepackage[a4paper,margin=1in]{geometry}
\usepackage{hyperref}
\usepackage{xurl}
\usepackage{booktabs}
\usepackage{amsmath}
\usepackage{amssymb}
\usepackage{listings}
\usepackage{xcolor}
\usepackage{enumitem}
\usepackage{parskip}

% Code listings (general-purpose).
\lstset{
  basicstyle=\ttfamily\small,
  keywordstyle=\color{blue},
  commentstyle=\color{gray},
  stringstyle=\color{teal},
  showstringspaces=false,
  columns=fullflexible,
  frame=single,
  framerule=0.2pt,
  breaklines=true
}

\setlist[itemize]{topsep=0.3em,itemsep=0.2em}
\setlist[enumerate]{topsep=0.3em,itemsep=0.2em}

% Simple per-section source line.
\newcommand{\notesources}[1]{\par\noindent{\small\color{gray}\textit{Sources:} \sloppy #1\par}}

% Unit-wise source strings for this subject (used by slides/notes).
% Path is relative to the lecture `latex/` folder (the compilation CWD).
\IfFileExists{../../../_shared/aiml-sources.tex}{%
  % Source strings for Elements of AIML (used by slides + notes).

\newcommand{\AIMLRepoURL}{https://github.com/tali7c/Elements-of-AIML}

\newcommand{\AIMLLabSources}{%
Provost and Fawcett, \textit{Data Science for Business}.;%
McKinney, \textit{Python for Data Analysis}.;%
WEKA Documentation (Waikato Environment for Knowledge Analysis), accessed 2026-02-28.;%
Matplotlib and Seaborn documentation, accessed 2026-02-28.%
}

\newcommand{\AIMLUnitISources}{%
Rich and Knight, \textit{Artificial Intelligence}, McGraw Hill, 2017.;%
Russell and Norvig, \textit{Artificial Intelligence: A Modern Approach} (4e), Ch. 1--2 (Introduction, Intelligent Agents).;%
Poole and Mackworth, \textit{Artificial Intelligence: Foundations of Computational Agents} (2e), selected sections.%
}

\newcommand{\AIMLUnitIISources}{%
Russell and Norvig, \textit{Artificial Intelligence: A Modern Approach} (4e), Ch. 7--9 (Logical Agents, First-Order Logic, Inference).;%
Huth and Ryan, \textit{Logic in Computer Science} (2e), selected sections on propositional and first-order logic.;%
Genesereth and Nilsson, \textit{Logical Foundations of Artificial Intelligence}, selected chapters.%
}

\newcommand{\AIMLUnitIIISources}{%
Mueller and Massaron, \textit{Machine Learning for Dummies}, For Dummies, 2016.;%
Mitchell, \textit{Machine Learning}, selected chapters on learning basics and evaluation.;%
Scikit-learn documentation, accessed 2026-02-28.%
}

\newcommand{\AIMLUnitIVSources}{%
Mueller and Massaron, \textit{Machine Learning for Dummies}, For Dummies, 2016.;%
James, Witten, Hastie, and Tibshirani, \textit{An Introduction to Statistical Learning} (2e), selected chapters.;%
Scikit-learn documentation, accessed 2026-02-28.%
}

\newcommand{\AIMLUnitVSources}{%
NIST, \textit{AI Risk Management Framework (AI RMF 1.0)}, 2023.;%
Barocas, Hardt, and Narayanan, \textit{Fairness and Machine Learning} (online book), selected sections.;%
Knaflic, \textit{Storytelling with Data}, selected chapters on communicating results.%
}
%
}{%
  % Faculty copies live under `private/` so their compile CWD is different.
  \IfFileExists{../../../../public/_shared/aiml-sources.tex}{%
    % Source strings for Elements of AIML (used by slides + notes).

\newcommand{\AIMLRepoURL}{https://github.com/tali7c/Elements-of-AIML}

\newcommand{\AIMLLabSources}{%
Provost and Fawcett, \textit{Data Science for Business}.;%
McKinney, \textit{Python for Data Analysis}.;%
WEKA Documentation (Waikato Environment for Knowledge Analysis), accessed 2026-02-28.;%
Matplotlib and Seaborn documentation, accessed 2026-02-28.%
}

\newcommand{\AIMLUnitISources}{%
Rich and Knight, \textit{Artificial Intelligence}, McGraw Hill, 2017.;%
Russell and Norvig, \textit{Artificial Intelligence: A Modern Approach} (4e), Ch. 1--2 (Introduction, Intelligent Agents).;%
Poole and Mackworth, \textit{Artificial Intelligence: Foundations of Computational Agents} (2e), selected sections.%
}

\newcommand{\AIMLUnitIISources}{%
Russell and Norvig, \textit{Artificial Intelligence: A Modern Approach} (4e), Ch. 7--9 (Logical Agents, First-Order Logic, Inference).;%
Huth and Ryan, \textit{Logic in Computer Science} (2e), selected sections on propositional and first-order logic.;%
Genesereth and Nilsson, \textit{Logical Foundations of Artificial Intelligence}, selected chapters.%
}

\newcommand{\AIMLUnitIIISources}{%
Mueller and Massaron, \textit{Machine Learning for Dummies}, For Dummies, 2016.;%
Mitchell, \textit{Machine Learning}, selected chapters on learning basics and evaluation.;%
Scikit-learn documentation, accessed 2026-02-28.%
}

\newcommand{\AIMLUnitIVSources}{%
Mueller and Massaron, \textit{Machine Learning for Dummies}, For Dummies, 2016.;%
James, Witten, Hastie, and Tibshirani, \textit{An Introduction to Statistical Learning} (2e), selected chapters.;%
Scikit-learn documentation, accessed 2026-02-28.%
}

\newcommand{\AIMLUnitVSources}{%
NIST, \textit{AI Risk Management Framework (AI RMF 1.0)}, 2023.;%
Barocas, Hardt, and Narayanan, \textit{Fairness and Machine Learning} (online book), selected sections.;%
Knaflic, \textit{Storytelling with Data}, selected chapters on communicating results.%
}
%
  }{%
    % Source strings for Elements of AIML (used by slides + notes).

\newcommand{\AIMLRepoURL}{https://github.com/tali7c/Elements-of-AIML}

\newcommand{\AIMLLabSources}{%
Provost and Fawcett, \textit{Data Science for Business}.;%
McKinney, \textit{Python for Data Analysis}.;%
WEKA Documentation (Waikato Environment for Knowledge Analysis), accessed 2026-02-28.;%
Matplotlib and Seaborn documentation, accessed 2026-02-28.%
}

\newcommand{\AIMLUnitISources}{%
Rich and Knight, \textit{Artificial Intelligence}, McGraw Hill, 2017.;%
Russell and Norvig, \textit{Artificial Intelligence: A Modern Approach} (4e), Ch. 1--2 (Introduction, Intelligent Agents).;%
Poole and Mackworth, \textit{Artificial Intelligence: Foundations of Computational Agents} (2e), selected sections.%
}

\newcommand{\AIMLUnitIISources}{%
Russell and Norvig, \textit{Artificial Intelligence: A Modern Approach} (4e), Ch. 7--9 (Logical Agents, First-Order Logic, Inference).;%
Huth and Ryan, \textit{Logic in Computer Science} (2e), selected sections on propositional and first-order logic.;%
Genesereth and Nilsson, \textit{Logical Foundations of Artificial Intelligence}, selected chapters.%
}

\newcommand{\AIMLUnitIIISources}{%
Mueller and Massaron, \textit{Machine Learning for Dummies}, For Dummies, 2016.;%
Mitchell, \textit{Machine Learning}, selected chapters on learning basics and evaluation.;%
Scikit-learn documentation, accessed 2026-02-28.%
}

\newcommand{\AIMLUnitIVSources}{%
Mueller and Massaron, \textit{Machine Learning for Dummies}, For Dummies, 2016.;%
James, Witten, Hastie, and Tibshirani, \textit{An Introduction to Statistical Learning} (2e), selected chapters.;%
Scikit-learn documentation, accessed 2026-02-28.%
}

\newcommand{\AIMLUnitVSources}{%
NIST, \textit{AI Risk Management Framework (AI RMF 1.0)}, 2023.;%
Barocas, Hardt, and Narayanan, \textit{Fairness and Machine Learning} (online book), selected sections.;%
Knaflic, \textit{Storytelling with Data}, selected chapters on communicating results.%
}
%
  }%
}


\title{Elements of AIML\\Unit 03 -- Lecture 03 Notes\\Dimensionality Reduction (PCA, LDA, ICA)}
\author{Tofik Ali}
\date{\today}

\begin{document}
\maketitle
\tableofcontents

\section{Lecture Overview}
This lecture introduces \textbf{dimensionality reduction} and focuses on PCA, with conceptual overviews of LDA and ICA.
Dimensionality reduction is important for:
\begin{itemize}[leftmargin=*]
  \item faster training/inference,
  \item reducing noise and redundancy,
  \item visualization of high-dimensional datasets.
\end{itemize}
\notesources{\AIMLUnitIIISources}

\section{How to use these notes (self-study)}
First, build intuition: PCA finds directions of maximum variance and projects data onto them.
Then study the step-by-step pipeline (centering, covariance, eigenvectors, projection) and work through the practice questions.
When using PCA in code, remember the evaluation rule: fit PCA on training data only to avoid leakage.

\section{Learning Outcomes}
\begin{itemize}[leftmargin=*]
  \item Explain the curse of dimensionality at an intuitive level.
  \item Describe PCA and interpret principal components and explained variance.
  \item Outline PCA steps and apply PCA using Python.
  \item Distinguish PCA vs LDA vs ICA conceptually.
\end{itemize}

\section{Keywords}
\begin{itemize}[leftmargin=*]
  \item \textbf{Dimensionality:} number of features ($d$).
  \item \textbf{Principal component:} direction capturing maximum variance, orthogonal to previous.
  \item \textbf{Explained variance:} variance captured by a component (often ratio).
  \item \textbf{Projection:} mapping data to a lower-dimensional space.
  \item \textbf{Covariance matrix:} captures relationships between features.
\end{itemize}

\section{Abbreviations and symbols}
\begin{itemize}[leftmargin=*]
  \item $n$: number of samples (rows), $d$: original number of features, $k$: reduced number of components.
  \item $X \in \mathbb{R}^{n \times d}$: data matrix; $X_c$: mean-centered data.
  \item $C$: covariance matrix; $\lambda_i$: eigenvalues; $v_i$: eigenvectors.
  \item $V_k$: matrix of top-$k$ eigenvectors; $Z = X_c V_k$: projected data.
  \item \textbf{EVR}: explained variance ratio $r_i = \lambda_i / \sum_j \lambda_j$.
\end{itemize}

\section{Why dimensionality reduction?}
\subsection{Curse of dimensionality (intuition)}
As dimensionality increases:
\begin{itemize}[leftmargin=*]
  \item data becomes sparse (needs more samples to cover space),
  \item distances can become less informative,
  \item models can overfit more easily.
\end{itemize}
Reducing dimensions can improve generalization and computational efficiency.

\subsection{When should you reduce dimensions?}
Common reasons:
\begin{itemize}[leftmargin=*]
  \item \textbf{Speed}: fewer features $\Rightarrow$ faster training and prediction.
  \item \textbf{Noise reduction}: remove low-variance directions that often contain noise.
  \item \textbf{Visualization}: project to 2D/3D for plots and intuition.
  \item \textbf{Redundancy}: features may be highly correlated; PCA compresses them.
\end{itemize}

\subsection{Dimensionality reduction vs feature selection}
\begin{itemize}[leftmargin=*]
  \item \textbf{Feature selection}: keep a subset of original features (easier to interpret).
  \item \textbf{Feature extraction}: create new features (components) from combinations of originals (e.g., PCA).
\end{itemize}

\section{Principal Component Analysis (PCA)}
\subsection{Idea}
PCA finds a new coordinate system where:
\begin{itemize}[leftmargin=*]
  \item the first axis captures the maximum variance,
  \item the second captures the next maximum variance while being orthogonal, etc.
\end{itemize}
Then we keep only the top $k$ axes and drop the rest.

\subsection{What is being optimized? (important intuition)}
PCA can be motivated in two equivalent ways:
\begin{itemize}[leftmargin=*]
  \item \textbf{Max-variance view:} choose a unit vector $v$ that maximizes the variance of the projection $X_c v$.
  \item \textbf{Min-reconstruction-error view:} choose a $k$-dimensional subspace that minimizes squared reconstruction error.
\end{itemize}
Both lead to the same principal components.

\subsection{Steps (more detail)}
Given data matrix $X$ with $n$ samples and $d$ features:
\begin{enumerate}[leftmargin=*]
  \item (Often) standardize features: subtract mean and divide by std.
  \item Mean-center: $X_c = X - \mu$ (where $\mu$ is the mean of each feature).
  \item Compute covariance matrix:
  \[
    C = \frac{1}{n-1} X_c^\top X_c
  \]
  \item Compute eigenvalues $\lambda_i$ and eigenvectors $v_i$ of $C$.
  \item Sort eigenvectors by decreasing eigenvalues.
  \item Choose top-$k$ eigenvectors and project: $Z = X_c V_k$.
\end{enumerate}

\subsection{Explained variance}
Eigenvalue $\lambda_i$ corresponds to variance captured by component $i$.
Explained variance ratio:
\[
  r_i = \frac{\lambda_i}{\sum_j \lambda_j}
\]
We often choose $k$ such that $\sum_{i=1}^k r_i$ exceeds a threshold (e.g., 0.90), depending on the application.

\subsection{How to choose $k$ in practice}
Three common approaches:
\begin{itemize}[leftmargin=*]
  \item \textbf{Cumulative EVR}: pick smallest $k$ reaching 90\%, 95\%, etc.
  \item \textbf{Scree plot (elbow)}: choose $k$ where additional components add little.
  \item \textbf{Downstream performance}: treat $k$ as a hyperparameter and validate it (careful: do this inside CV to avoid leakage).
\end{itemize}

\subsection{Why standardize?}
If features have different scales (e.g., income in thousands vs age in years), PCA may focus on the largest-scale feature. Standardization makes PCA reflect correlation structure rather than raw scale.

\subsection{A tiny geometric example (2D $\rightarrow$ 1D)}
Suppose (after centering) the points lie roughly on the line $x_2 \approx x_1$.
Then the first principal component is approximately:
\[
  v_1 = \frac{1}{\sqrt{2}}\begin{bmatrix}1\\1\end{bmatrix}
\]
Projecting a point $x$ to 1D gives the scalar coordinate:
\[
  z = x^\top v_1
\]
Intuition: we keep the ``diagonal direction'' and drop the perpendicular direction, losing little information if points are tightly clustered around the diagonal.

\subsection{PCA and SVD (brief note)}
Many software libraries compute PCA using \textbf{SVD} (singular value decomposition) for numerical stability.
You do not need to compute SVD by hand, but it explains why PCA works well for large matrices.

\subsection{Common pitfalls}
\begin{itemize}[leftmargin=*]
  \item PCA is \textbf{linear}: it cannot capture complex non-linear structure (for that, methods like t-SNE/UMAP exist).
  \item Components can be hard to interpret; they are combinations of features.
  \item Always fit PCA on \textbf{training data only} (otherwise you leak information from validation/test).
\end{itemize}

\section{LDA and ICA (conceptual)}
\subsection{LDA}
Linear Discriminant Analysis is supervised: it uses class labels and tries to find a projection that separates classes.
\paragraph{How it differs from PCA.}
PCA ignores labels and maximizes variance overall.
LDA uses labels and tries to maximize between-class separation while minimizing within-class spread.

\subsection{ICA}
Independent Component Analysis tries to find components that are statistically independent. A classic intuition is separating mixed signals into original sources.
\paragraph{Where ICA is used.}
Blind source separation (signals), artifact removal, and some feature extraction tasks.
ICA assumes the underlying sources are independent and non-Gaussian.

\section{Python example (PCA)}
\begin{lstlisting}[language=Python]
from sklearn.decomposition import PCA
from sklearn.preprocessing import StandardScaler

X_scaled = StandardScaler().fit_transform(X)
pca = PCA(n_components=2)
Z = pca.fit_transform(X_scaled)
print(pca.explained_variance_ratio_)
\end{lstlisting}

\section{Practice (with answers)}
\subsection*{P1. Why is it incorrect to apply PCA on the full dataset before splitting into train/test?}
\textbf{Answer:} PCA learns directions from data. If you fit PCA using test data, information from test influences the representation used in training, causing leakage and optimistic evaluation.

\subsection*{P2. PCA vs feature selection: which is easier to interpret and why?}
\textbf{Answer:} Feature selection is easier to interpret because it keeps original features. PCA creates components that mix features, which can be harder to explain.

\subsection*{P3. If EVR values are $[0.60, 0.25, 0.10, 0.05]$, what is cumulative EVR for $k=2$?}
\textbf{Answer:} $0.60 + 0.25 = 0.85$ (85\%).

\section{Exit questions (with answers)}

\subsection*{Q1. What does PCA optimize?}
\textbf{Answer:} PCA chooses orthogonal directions that maximize variance of the projected data (sequentially), thus capturing as much variability as possible in fewer dimensions.

\subsection*{Q2. Why standardize before PCA?}
\textbf{Answer:} to avoid features with larger numeric scales dominating the covariance and the principal components.

\subsection*{Q3. Compare PCA and LDA.}
\textbf{Answer:} PCA is unsupervised and maximizes variance; LDA is supervised and aims to separate classes by maximizing between-class separation relative to within-class scatter.

\subsection*{Q4. What is explained variance?}
\textbf{Answer:} the amount (or proportion) of total variance in the data captured by a principal component; used to choose the number of components.

\section{Summary}
PCA is a core dimensionality reduction tool. LDA and ICA provide alternative projections depending on whether labels are available (LDA) or independence is desired (ICA).
\notesources{\AIMLUnitIIISources}

\end{document}
